\documentclass[12pt]{article}

\usepackage{graphicx, color}
\usepackage{setspace}  %espaçamento entre linhas ABNT
\usepackage{indentfirst}
\usepackage{color}
\usepackage{epigraph} %para quem quer colocar epigrafe
\usepackage{todonotes}
%\parindent=15ex
%\oddsidemargin=1cm
%\topmargin=-0.30cm
%\textheight=22.5cm
%\textwidth=15.1cm
%\def \baselinestretch{1.5}
\usepackage{scalefnt} %para mudar o tamanho da fonte
\usepackage{amsmath}
%Pacote Indice remissivo
\usepackage{makeidx}
\usepackage[english]{babel} % pacotes para figura/tabelas...
\usepackage{graphicx, color}
\usepackage{threeparttable}
\usepackage{epigraph} %para quem quer colocar epigrafe
\usepackage[all]{xy}
\usepackage[utf8]{inputenc}
\usepackage{amssymb}
\usepackage{a4wide}
\usepackage{aecompl}
\usepackage{exscale}
% \usepackage{cite}
\usepackage{adjustbox}
\usepackage{booktabs}
%\usepackage[num]{abntcite}
\usepackage[final]{hyperref} % adds hyper links inside the generated pdf file
\usepackage{listings}
\usepackage{titlesec}
\usepackage{biblatex}
\usepackage{rotating}
\usepackage{comment}
\usepackage{svg}
\usepackage{color, colortbl}


\hypersetup{
	colorlinks=true,       % false: boxed links; true: colored links
	linkcolor=blue,        % color of internal links
	citecolor=blue,        % color of links to bibliography
	filecolor=magenta,     % color of file links
	urlcolor=blue         
}

%formatacao da section
\titleformat{\section}    
       {\normalfont\fontfamily{phv}\fontsize{16}{17}\bfseries}
       {\thesection.}{0.4em}{}
\titleformat{\subsection}    
       {\normalfont\fontfamily{phv}\fontsize{14}{17}\bfseries}
       {\thesubsection.}{0.1em}{}
       
\titlelabel{\thetitle.\quad}
\titlespacing{\section}{2pt}{20pt}{8pt}
%\titlelabel{\thetitle.\bfseries\quad} %coloca o "." depois.


\pagenumbering{arabic}

\let\oldabstract\abstract
\let\oldendabstract\endabstract
\makeatletter
\renewenvironment{abstract}
{\renewenvironment{quotation}%
               {\list{}{\addtolength{\leftmargin}{0.1cm}
                        \listparindent 0cm%
                        \itemindent    \listparindent%
                        \rightmargin   \leftmargin%
                        \parsep        \z@ \@plus\p@}%
                \item\relax}%
               {\endlist}%
\oldabstract}
{\oldendabstract}
\makeatother

\addbibresource{proposta.bib}

\interfootnotelinepenalty=10000


\graphicspath{{images/}}
\usepackage{subcaption}


\begin{document}
\onehalfspacing
%%%%%%%%%%%%%%%%%%%%%%%%%%%%%%%%% CAPA %%%%%%%%%%%%%%%%%%%%%%%%%%%%%%%%%
%\pagestyle{empty}

\begin{center}
University of São Paulo\\
Institute of Mathematics and Statistics\\
Computer Science Department\\
--\\
\textit{Research Project for Undergraduate Research
\\São Paulo Research Foundation (FAPESP)}\\

\hfill\\

\begin{Large}
    \textbf{An analysis of the delay Linux patches experience from their upstream inclusion to LTS release} %TEMPORARY
    %
    %    - Between development and release: An analysis of the delay between patch integration and LTS inclusion
    %    - Between development and release of Linux patches: An analysis of the delay patches experience before being ported to LTS releases
    %    - An analysis of the delay Linux patches experience from their upstream inclusion to LTS release
    %    - (Ideia David) An analysis of backporting lag in Linux LTS releases
    %
\end{Large}

\hfill\\

\textbf{Student}: Vitor Gomes da Silva Messias\\
\textbf{Advisors}: Paulo Meirelles, David Tadokoro\\
\textit{vitorgdsmessias@usp.br, paulormm@ime.usp.br, davidbtadokoro@usp.br}
\end{center}


\begin{abstract}
The Linux Kernel, a FLOSS project with over 30 years of active development, is a cornerstone of modern computing. Along with that, it is also a massive project that involves thousands of developers from across the globe and is always in constant change. However, its widespread usage means many hard to update systems cannot keep up with its changes and need more stable versions of the kernel. As such, the project provides many Long Term Support (LTS) versions, these are actively maintained versions that strive for maintainability. Consequently, maintaining older versions of a software means that important bugfixes made to the main Linux repository (upstream) need to be ported over to those versions, this process is called backporting. This is the process this research project aims to analyse. More specifically, it aims to provide an analysis of the time in which patches (changes) to the main Linux Repository take to be backported over to LTS branches. Since many modern computing projects rely on LTS versions, it stands to reason that critical changes to upstream, such as security updates, must be quickly ported to LTS. With this analysis the research aims to 1) provide a visualisation of data, to give developers quick insight in bottlenecks regarding porting, and 2) provide tools or solutions to such bottlenecks. 

\textbf{Keywords}: Linux, FLOSS, Linux kernel development model, FLOSS development model, bottleneck mitigation, software engineering, backporting, LTS, lag
visualisation.

\end{abstract}

\vfill

\section{Introduction}
\label{chp:introduction} 

The Operating System(OS) is an integral part of the vast majority of computing devices, with the kernel being it's core component. As such, the Linux kernel fills an important role in the computing world, being a Free/Libre/Open Source Software (FLOSS) that's highly customizable and free of licensing costs \cite{sbes}. Because of that, it is widely used across many sectors, such as machine learning, and embedded systems. 

\subsection{The Linux kernel project structure}
    The Linux kernel is a piece of software thats actively developed by thousands of developers across the globe, seeing contributions even from some companies. To properly organize such a huge amount of contributors the kernel is divided into subsystems, with each one of those having dedicated maintainers who oversee contributions to them.\cite{sbes}

    Those contributions are made by patches, which are text documents that describe the diference between two source trees. But contributions to the subsystems do not go directly to the Linus Torvald's repository(Linux's main tree, also called the upstream). Instead each one of the subsystems maintains a separate source tree and, at the start of development cycles selected changes to these trees are merged into the upstream. \cite{KernelDocs}

\subsubsection{The Linux kernel Versioning System}

    To accomodate for the constantly changing upstream the kernel has a robust versioning system. Before the current version is considered stable enough for release, the kernel has intermediate versions appended with rc, these are meant for developers and linux entusiasts and may break things. After this period, the next stable release is then launched, this is the ready for use version, and it aims to not have caused any regressions, which means that it aims to not break any existing systems in the old version.
    
    After the stable release a new development cycle starts for the next version. However, while development to the next version may have already started, the stable release is still updated with bugfixes, and these are made at least until the next stable version releases. \cite{KernelDocs}

    Moreover, since Linux is used in many critical and hard to update systems, the kernel also  provides Long Term Support (LTS) versions. These receive only critical bugfixes and security updates for longer periods of time (usually from 2 to 6). \cite{KernelReleases}

    \begin{figure}[htbp]
    \centering
        \includegraphics[width=0.5\textwidth]{LinuxVersioning.pdf}
    \caption{Lifecycle of Linux kernel versions \cite{KernelCycle}}
    \end{figure}

\subsubsection{Backporting}

    The Linux kernel versioning and release system means that it has many supported stable versions that rely on older development trees. This raises the issue of getting patches, which are usually bugfixes and security concerns, implemented into these stable versions, as the bugfixes implemented upstream may depend on features that are only on future versions.

    The task of adapting patches from newer versions of the software, to be compatible with older versions of software is called backporting. And this work is essential to properly maintain the kernel's stable releases.

\subsection{Problem Outline}

    As LTS versions are used across many critical infrastructures across the computing world, it becomes essencial that those versions are free of major bugs, and more importantly, security issues. As such, patches that address these issues upstream must be quickly backported to all LTS versions.

    This is because the time between a major security issue being fixed upstream, and it being ported into the LTS versions leaves a gap of vulnerability. And, during this gap many systems are vulnerable to known exploits.

    This however poses a major challenge, as one patch has to be adapted to several different LTS versions, each one of which needing differnet adaptations. That means there exists an inherent delay associated with the backporting of patches.\cite{icpc22-backporting}

    So while we cannot nullify the delay, it is still important to keep it to a mininum. Therefore, reasearching more in-depth how this backporting delay presents itself on the Linux Project can aid decisions in the current and future of the project's development.
    
\section{Objectives}
\label{chp:objectives}

Sketch of Objectives

1. Analyse and understand the lag in backporting to LTS.

2. Provide a visualisation so that project participants can take action

3. Provide potential insight into potential pain points in backporting
\section{Justification}
\label{chp:justification}

Sketch of Justifications

1. Outline the merits and experience our research team has about reasearching the kernel.

2. Outline all the contributions our team has done to the Linux kernel.

(both 1. and 2. are probably going to be copies of David's proposal)

3.Outline the importance LTS patches have to modern systems, and how they could be severely affected by delays(especially when talking about security patches)

4. Give indicators that we might have a bottleneck in backporting on some areas of the kernel (dificult considering the original paper hasn't been published).

5. síntese da bibliografia fundamental as written on \href{FAPESP's page}{https://fapesp.br/253/projeto-de-pesquisa}
\section{Research Methods}
\label{chp:research-methods}

TBD
\section{Work Plan and Execution Schedule}
\label{chp:workplan-schedule}

Por enquanto vou deixar aqui o cronograma feito para MAC0215:
\subsection{cronograma}

- Agosto:\\
    - Leitura de artigos indicados pelo orientador (10hrs)\\
    - Escrita de proposta inicial de IC FAPESP (10hrs)\\
- Setembro:\\
    - Revisão e finalização da proposta de IC FAPESP (10hrs)\\
    - Aclimatação com o projeto kworkflow e setup do ambiente de desenvolvimento (10hrs)\\
    - Contribuições iniciais ao projeto kworkflow (15hrs)\\
- Outubro:\\
    - Contribuições avançadas ao projeto kworkflow (20hrs)\\
- Novembro:\\
    - Planejamento de artigo relacionado à IC FAPESP (15hrs)\\
    - Escrita, revisão e submissão de artigo relacionado à IC FAPESP \\(15hrs)

\section{Publications}
\label{sec:publications}

Acho que não será aplicável?

\printbibliography
\end{document}
